\documentclass[11pt]{article}

\usepackage[T1]{fontenc}
\usepackage[utf8]{inputenc}
\usepackage{lmodern}
\usepackage{microtype}
\usepackage[margin=0.95in]{geometry}
\usepackage{amsmath,amssymb,amsthm,mathtools}
\usepackage{array,booktabs,longtable}
\usepackage{xcolor}
\usepackage[colorlinks=true,linkcolor=blue!55!black,urlcolor=blue!55!black]{hyperref}

\setlength{\emergencystretch}{3em}
\setlength{\parskip}{0.45em}
\setlength{\parindent}{0pt}

\newtheorem{theorem}{Theorem}[section]
\newtheorem{proposition}[theorem]{Proposition}
\newtheorem{lemma}[theorem]{Lemma}
\newtheorem{corollary}[theorem]{Corollary}
\theoremstyle{remark}
\newtheorem{remark}[theorem]{Remark}

\newcommand{\R}{\mathbb R}
\newcommand{\N}{\mathbb N}
\newcommand{\E}{\mathbb E}
\newcommand{\Normal}{\mathcal N}
\newcommand{\TV}{\operatorname{TV}}
\newcommand{\AO}{\mathrm{AO}}
\newcolumntype{L}[1]{>{\raggedright\arraybackslash}p{#1}}

\title{Full-Batch Optimality at Fixed Participation\\
A Comparison of Transcript and Average Privacy}
\author{}
\date{}

\begin{document}
\maketitle

\begin{abstract}
We compare the variance of unbiased frozen-gradient averages when each record
has a fixed expected participation budget $E\in\N_{>0}$. A run of $T\geq E$ releases
uses independent Poisson sampling at rate $E/T$. For Gaussian noise calibrated
to a common $(\varepsilon,\delta)$ target, let $v_{\mathrm{FB}}$ be the
full-batch variance and $a=\Delta/\sqrt{v_{\mathrm{FB}}}$.
If $\varepsilon>0$ and
$a^2\leq6E\varepsilon/(3E+2)$, full batch has strictly smaller variance
than every smaller batch, at every nonnegative frozen-gradient energy.
The conclusion holds whether privacy protects the canonical average alone
or the entire canonical transcript. We prove a sharper horizon-dependent
condition and give a positive-noise counterexample to unconditional optimality.
For transcript calibration we also determine the exact sparse critical-window
profile and calibrated signal-to-noise offset. These imply eventual full-batch
dominance when $0<\delta<1-e^{-E}$. In contrast, calibration of the average
alone admits a Gaussian variance bounded uniformly in the horizon.
\end{abstract}

\section{The comparison}
Fixing how often a record is expected to participate leaves a choice: include
it in a few large batches, or sample it less often across more releases.
The average remains unbiased in either case. Its variance has two terms,
one from participation and one from the Gaussian noise needed for privacy.
Our question is whether full batch minimizes their sum.

A tail event of the average gives an explicit sufficient condition for strict
full-batch optimality at every finite horizon. The same argument works when
privacy protects the average alone or the complete transcript. An adjacent
smaller-batch counterexample shows why some condition is necessary.
The sparse limit then separates the two output choices: the complete transcript
reveals individual large observations, while the average retains only their sum.

We derive the variance and calibration in Section~\ref{sec:model}, prove the
finite comparison in Section~\ref{sec:near-full}, and analyze the sparse limit
in Section~\ref{sec:sparse}. The model is an equal-weight average of a fixed
recordwise query; it makes no assertion about adaptive optimizer trajectories
or training loss.

\section{A fixed budget and two public outputs}\label{sec:model}
Fix a positive integer $E$, a clipping bound $\Delta>0$, and
$\varepsilon\geq0$, $0<\delta<1$. For a dataset of size $N$, hold the query
fixed and write its contributions as $\widetilde g_1,\ldots,\widetilde g_N$,
with $|\widetilde g_i|\leq\Delta$. Define
\[
 G=\sum_i\widetilde g_i,\qquad S=\sum_i\widetilde g_i^2.
\]
For an integer $T\geq E$, set $q=E/T$. Independent indicators
$B_{it}\sim\operatorname{Bernoulli}(q)$ and independent noises
$\xi_t\sim\Normal(0,\sigma^2)$ give the estimator
\[
 \widehat G_{E,T}=\frac1E\sum_{t=1}^T
       \left(\sum_i B_{it}\widetilde g_i+\xi_t\right).
\]
Since $Tq=E$, its expectation is $G$. Independence gives the variance directly:
\begin{equation}\label{eq:variance}
 \operatorname{Var}(\widehat G_{E,T})
 =\frac{Tq(1-q)S+T\sigma^2}{E^2}
 =\frac{1-q}{E}S+v,\qquad v:=\frac{T\sigma^2}{E^2}.
\end{equation}
Thus $v$ is the Gaussian variance of the average, whereas $\sigma^2$ is the
noise variance of one raw release. We use full-sum normalization throughout.
The expected batch size is $qN$; $q$ itself is held fixed across neighboring
datasets. Full batch means $T=E$.

\paragraph{The canonical pair.}
Under add/remove adjacency the maximal differing contribution is $\Delta$.
At raw noise standard deviation $\sigma$, the canonical one-release laws are
\[
 P_{q,\sigma}=(1-q)\Normal(0,\sigma^2)+q\Normal(\Delta,\sigma^2),
 \qquad Q_\sigma=\Normal(0,\sigma^2).
\]
Here and below a zero-variance Gaussian is a point mass. For arbitrary laws,
put
\[
 H_\alpha(P,Q)=\sup_A\{P(A)-\alpha Q(A)\},\qquad
 \mathcal H_\varepsilon(P,Q)
 =\max\{H_{e^\varepsilon}(P,Q),H_{e^\varepsilon}(Q,P)\}.
\]
The supremum is over measurable events. For use with a general schedule write
\[
 \AO_S(\sigma,\Delta,q,T,\varepsilon)
 :=\mathcal H_\varepsilon(P_{q,\sigma}^{\otimes T},Q_\sigma^{\otimes T});
\]
the subscript denotes the sampled-Gaussian experiment, not the energy $S$.

At a supplied variance $v$, the transcript uses raw noise
$\sigma=E\sqrt{v/T}$. Its average has the pair
\begin{equation}\label{eq:average-experiment}
 \overline P_{E,T,v}=\E_K\Normal(K\Delta/E,v),\qquad
 \overline Q_v=\Normal(0,v),\qquad
 K\sim\operatorname{Binomial}(T,E/T),
\end{equation}
where $\E_K\Normal(K\Delta/E,v)$ denotes a mixture. Define the two profiles
and their calibrated variances by
\begin{align*}
 \mathcal A_{E,T}(v;\varepsilon)
 &=\AO_S(E\sqrt{v/T},\Delta,E/T,T,\varepsilon),&
 v_{E,T}&=\inf\{v\geq0:\mathcal A_{E,T}(v;\varepsilon)\leq\delta\},\\
 \overline{\mathcal A}_{E,T}(v;\varepsilon)
 &=\mathcal H_\varepsilon(\overline P_{E,T,v},\overline Q_v),&
 \overline v_{E,T}&=\inf\{v\geq0:\overline{\mathcal A}_{E,T}(v;\varepsilon)\leq\delta\}.
\end{align*}
Averaging is a common channel on the two transcript laws. Data processing gives
\begin{equation}\label{eq:average-transcript-order}
 \overline{\mathcal A}_{E,T}(v;\varepsilon)\leq\mathcal A_{E,T}(v;\varepsilon),
 \qquad \overline v_{E,T}\leq v_{E,T}.
\end{equation}
Set $\sigma_{E,T}=E\sqrt{v_{E,T}/T}$ and
\[
 V_E(T,S)=\frac{1-q}{E}S+v_{E,T},\qquad
 \overline V_E(T,S)=\frac{1-q}{E}S+\overline v_{E,T}.
\]
These calibrations protect different public outputs. Average-only calibration
provides no guarantee for additionally disclosed individual releases.

The canonical profile bounds the privacy profile of every neighboring pair
for the fixed clipped query. To see the mechanism, a contribution
$|d|\leq\Delta$ is obtained by scaling the canonical observation by $d/\Delta$
and adding independent Gaussian noise to restore its variance. Adding the
shared sampled contributions is another common channel. Appendix~\ref{app:canonical-domination}
gives both constructions explicitly. Noise is calibrated uniformly over the
query class; $S$ describes the particular dataset used for the variance comparison.

Write $\phi,\Phi$ for the standard Gaussian density and CDF, and
$\overline\Phi=1-\Phi$.
\begin{lemma}[Calibration]\label{lem:calibration}
For each finite admissible horizon, both profiles are continuous and
nonincreasing in $v\geq0$. Their feasible sets are nonempty and their minima
are attained. At a positive minimum the profile equals $\delta$.
At full batch the two calibrated variances agree at a positive value
$v_{\mathrm{FB}}$ independent of $E$. Writing $a=\Delta/\sqrt{v_{\mathrm{FB}}}$,
\begin{equation}\label{eq:gaussian-calibration}
 \delta=G_\varepsilon(a):=
 \Phi\!\left(\frac a2-\frac\varepsilon a\right)
 -e^\varepsilon\Phi\!\left(-\frac a2-\frac\varepsilon a\right),\qquad
 G_\varepsilon'(a)=\phi(a/2-\varepsilon/a)>0.
\end{equation}
\end{lemma}
Appendix~\ref{app:calibration} proves the lemma,
including the singular zero-noise endpoint. Formula~\eqref{eq:gaussian-calibration}
is the analytical Gaussian calibration~\cite{ballewang2018}.
A consequence we will use repeatedly is
\begin{equation}\label{eq:calibration-transfer}
 \mathcal A_{E,T}(v_*;\varepsilon)>\delta\ \Longrightarrow\ v_{E,T}>v_*,
 \qquad
 \overline{\mathcal A}_{E,T}(v_*;\varepsilon)>\delta\ \Longrightarrow\ \overline v_{E,T}>v_*.
\end{equation}
Indeed, an attained feasible minimum at or below $v_*$ would make $v_*$
feasible by monotonicity. At full batch the sampling term vanishes, so both
total variances equal $v_{\mathrm{FB}}$.

\section{The finite comparison}\label{sec:near-full}
Full batch has no sampling variance. To prove that it wins at every energy,
we therefore compare only the calibrated Gaussian variances. The useful
order is
\[
 v_{\mathrm{FB}}<\overline v_{E,T}\leq v_{E,T}.
\]
The second inequality is data processing. For the first, we must find an
event of the average whose privacy value exceeds $\delta$ when its Gaussian
variance is still $v_{\mathrm{FB}}$.

\paragraph{The event and its count gap.}
Write $a=\Delta/\sqrt{v_{\mathrm{FB}}}$. At this variance the standardized
average has reference law $\Normal(0,1)$ and source law
$\E_K\Normal(aK/E,1)$. At full batch $K=E$, and the likelihood ratio
$e^{ay-a^2/2}$ exceeds $e^\varepsilon$ on the ray
$y>c_*:=a/2+\varepsilon/a$. Its privacy value is
\begin{equation}\label{eq:fullbatch-event-identity}
 \delta=\overline\Phi(c_*-a)-e^\varepsilon\overline\Phi(c_*).
\end{equation}
For a smaller batch, set
\[
 Z=K-E,\qquad h=a/E,\qquad
 \theta=E(\varepsilon/a^2-1/2).
\]
Then $aK/E-c_*=h(Z-\theta)$ and $a-c_*=-h\theta$. The reference term
in the event value is unchanged. Its entire change is consequently
\begin{equation}\label{eq:main-count-gap}
 \Gamma_{E,T}(\theta,h)=\E\Phi(h(Z-\theta))-\Phi(-h\theta),
\end{equation}
which gives
\begin{equation}\label{eq:event-profile-bound}
 \overline{\mathcal A}_{E,T}(v_{\mathrm{FB}};\varepsilon)
 \geq\delta+\Gamma_{E,T}(\theta,h).
\end{equation}
The center $\theta$ measures how far the full-batch threshold lies above its
signal mean, in units of one participation. Establishing a positive gap is
the substantive step: the Gaussian CDF is not convex on the whole real line.

\paragraph{Where the condition comes from.}
The count has $\E Z=0$, $\E Z^2=E(1-q)$ and
$\E Z^3=E(1-q)(1-2q)$. Thus, for fixed $E,T,\theta$, expansion at $h=0$ gives
\begin{equation}\label{eq:gap-small-scale}
 \Gamma_{E,T}(\theta,h)
 =\frac{\phi(0)\operatorname{Var}(Z)}2
   \left(\theta-\frac{1-2E/T}{3}\right)h^3+O(h^5).
\end{equation}
For low sampling rates this identifies a necessary center for positivity
at every scale. Theorem~\ref{thm:rate-adaptive-center} proves its sufficiency:
\[
 \Gamma_{E,T}(\theta,h)>0\quad\text{if}\quad
 \begin{cases}
 \theta\geq(1-2E/T)/3,&T>2E,\\
 \theta>0,&T=2E,\\
 \theta\geq0,&E<T<2E.
 \end{cases}
\]
At the low-rate boundary the cubic term itself vanishes. The appendix
establishes the sign at every positive scale by writing the scale derivative
as a Laplace integral of a nonnegative hinge transform. A binomial telescoping
identity supplies the required nonnegativity.

Solving these center inequalities for $a^2$ leads to
\begin{equation}\label{eq:rate-adaptive-signal-factor}
 \Lambda_{E,T}=
 \begin{cases}6ET/(3ET+2T-4E),&T>2E,\\2,&E<T\leq2E.\end{cases}
\end{equation}
The zero gap at the symmetric boundary will require a slightly different event.

\begin{theorem}[Finite-horizon comparison]\label{thm:rate-adaptive-pointwise}
Let $T>E>0$ be integers, $\Delta>0$, $\varepsilon>0$, and $0<\delta<1$.
If $a=\Delta/\sqrt{v_{\mathrm{FB}}}$ satisfies
$a^2\leq\Lambda_{E,T}\varepsilon$, then
\[
 v_{\mathrm{FB}}<\overline v_{E,T}\leq v_{E,T},\qquad
 V_E(E,S)=\overline V_E(E,S)<\overline V_E(T,S)\leq V_E(T,S)
 \quad(S\geq0).
\]
\end{theorem}
\begin{proof}
The preceding reduction and center comparison make the average-only profile
strictly larger than $\delta$, except possibly at $T=2E$, $a^2=2\varepsilon$.
Indeed, for $T>2E$,
\[
 \theta\geq\frac{1-2E/T}{3}
 \iff\frac{\varepsilon}{a^2}\geq\frac12+\frac{1-2E/T}{3E}
 =\Lambda_{E,T}^{-1};
\]
for the other rates, $\theta\geq0$ is equivalent to $a^2\leq2\varepsilon$.

At the remaining equality, $c_*=a$. Symmetry makes the ray $(a,\infty)$
tie the full-batch value, but the average's mixture likelihood ratio satisfies
\[
 L(a)=e^{a^2/2}\E\exp[-a^2(K/E-1)^2/2]<e^{a^2/2}=e^\varepsilon.
\]
A small interval just above $a$ therefore contributes negatively to that ray's
privacy value. Removing it gives a value strictly above $\delta$, as shown
in Lemma~\ref{lem:strict-average-endpoint}. This uses only the average.

In all cases~\eqref{eq:calibration-transfer} gives
$\overline v_{E,T}>v_{\mathrm{FB}}$. Apply data processing and add the
nonnegative sampling variance to finish.
\end{proof}

\begin{corollary}[All finite horizons]\label{cor:fixed-epoch-all-horizon}
Fix $E\in\N_{>0}$, $\Delta>0$, $\varepsilon>0$, and $0<\delta<1$.
If $a=\Delta/\sqrt{v_{\mathrm{FB}}}$ satisfies
\begin{equation}\label{eq:epoch-uniform-condition}
 a^2\leq\Lambda_E\varepsilon,\qquad \Lambda_E=\frac{6E}{3E+2},
\end{equation}
then for every integer $T>E$,
\[
 v_{\mathrm{FB}}<\overline v_{E,T}\leq v_{E,T},
\]
and, simultaneously for every $S\geq0$,
\[
 V_E(E,S)=\overline V_E(E,S)<\overline V_E(T,S)\leq V_E(T,S).
\]
\end{corollary}
\begin{proof}
For $T>2E$, $\Lambda_{E,T}=6E/(3E+2-4E/T)>\Lambda_E$.
For $E<T\leq2E$, $\Lambda_{E,T}=2>\Lambda_E$.
Apply Theorem~\ref{thm:rate-adaptive-pointwise}.
\end{proof}
The uniform hypothesis is exactly $\theta\geq1/3$. Every finite low-rate
boundary $(1-2E/T)/3$ is strictly below $1/3$, which explains why the closed
uniform condition still gives strict dominance. Equivalently, strict
monotonicity of $G_\varepsilon$ expresses the condition as
$\delta\leq G_\varepsilon(\sqrt{\Lambda_E\varepsilon})$.
The small-scale expansion proves sharpness of the low-rate center for this
particular event uniformly over $h$; it does not classify the complete
privacy profile outside the sufficient domain.

\section{A positive-noise reversal}
A qualification on privacy is necessary. At zero noise both canonical
profiles equal $1-(1-E/T)^T$, tending to $1-e^{-E}$. The next example lies
below this limiting boundary and has positive calibrated noise at both horizons.
\begin{theorem}[An adjacent smaller batch can win]\label{thm:strict-interior-counterexample}
For $\Delta=1$, $E=2$, $\varepsilon=1/4$, and $\delta=3/4$, write
$v_T=v_{2,T}$. Then
\[
 v_3\leq\frac5{32}<\frac4{25}<v_2,
 \qquad
 V_2(3,S)<V_2(2,S)\quad\text{for }0\leq S\leq\frac9{400}.
\]
The same strict total-variance comparison holds for average-only calibration.
\end{theorem}
\begin{proof}
Appendix~\ref{app:rational-calculation} bounds the forward and reverse
transcript profiles at $T=3$, $v=5/32$ by
$7420364678/10^{10}$ and $7237647581/10^{10}$, respectively. Both are below
$3/4$. At full batch, $v=4/25$ gives
$G_{1/4}(5/2)>19/25>3/4$. Feasibility and~\eqref{eq:calibration-transfer}
give the claimed brackets. For the total variance,
\[
 V_2(3,S)\leq\frac S6+\frac5{32}\leq\frac4{25}<v_2=V_2(2,S),
 \qquad 6(4/25-5/32)=9/400.
\]
Average-only calibration can only reduce the smaller-batch variance and
agrees with transcript calibration at full batch.
\end{proof}
The zero-noise profiles are $1$ and $26/27$, both larger than $3/4$.
Continuity therefore makes both calibrated noises positive. Also
$e^2>1+2+2>4$, so $3/4<1-e^{-2}$. The example disproves unconditional
optimality even below the nonparticipation boundary. Its loose privacy target
does not describe the prevalence of reversals at stringent targets.

\section{What the sparse limit reveals}\label{sec:sparse}
Now hold $E$ fixed and let $T\to\infty$. The count of participations converges
to a Poisson variable of mean $E$. Selection becomes rare per release, but
the probability of any selection does not vanish. Whether individual releases
are public is therefore consequential.

For transcript calibration, consider a raw-noise sequence with standardized
shift $\Delta/\sigma_T=\sqrt{2\log T}+s+o(1)$. A maximum test at threshold
$\sigma_T\sqrt{2\log T}$ has reference probability at most
$T\overline\Phi(\sqrt{2\log T})\to0$. A selected coordinate exceeds it with
probability tending to $\Phi(s)$, and the number of such exceedances tends
to Poisson$(E\Phi(s))$. This gives the lower bound $1-e^{-E\Phi(s)}$.
The likelihood-ratio argument in Appendix~\ref{app:sparse-limit} proves that
no event improves its limiting value.

\begin{theorem}[Sparse transcript calibration]\label{thm:smallbatch}
Fix $E\in\N_{>0}$, $\Delta>0$, $\varepsilon\geq0$, and
\begin{equation}\label{eq:interior}
 0<\delta<1-e^{-E}.
\end{equation}
Then, along integer horizons,
\[
 \frac\Delta{\sigma_{E,T}}-\sqrt{2\log T}
 \longrightarrow s_{E,\delta}:=
 \Phi^{-1}\!\left(-\frac{\log(1-\delta)}E\right).
\]
Consequently,
\begin{align}
 \sigma_{E,T}&\sim\frac\Delta{\sqrt{2\log T}}\longrightarrow0,
 \label{eq:sigmaasymptotic}\\
 v_{E,T}&\sim\frac{\Delta^2T}{2E^2\log T}\longrightarrow\infty,
 \label{eq:privacydivergence}\\
 V_E(T,S)&\sim v_{E,T}\qquad\text{for each fixed }S\geq0.
 \label{eq:total-equivalent}
\end{align}
\end{theorem}
The exact profile limit is inverted by bracketing the calibrated noise between
$\Delta/(\sqrt{2\log T}+s_{E,\delta}\pm\eta)$ for each $\eta>0$.
The general profile limit and this inversion are proved in the appendix.
The privacy variance alone eventually exceeds $v_{\mathrm{FB}}$, giving:
\begin{corollary}[One cutoff for all energies]\label{cor:eventual}
Under the hypotheses of Theorem~\ref{thm:smallbatch}, some integer $T_0>E$
satisfies $V_E(E,S)<V_E(T,S)$ for every $T\geq T_0$ and every $S\geq0$.
\end{corollary}
\begin{proof}
Choose $T_0$ with $v_{E,T}>v_{\mathrm{FB}}$ for $T\geq T_0$, and use
$(1-q)S/E\geq0$. The choice does not depend on $S$.
\end{proof}

\begin{proposition}[Average-only calibration stays bounded]\label{prop:average-uniform-bound}
For integers $T\geq E>0$, $\Delta>0$, $\varepsilon\geq0$, and $0<\delta<1$,
\[
 \overline v_{E,T}\leq\frac{\Delta^2}{2\pi\delta^2},\qquad
 \overline V_E(T,S)\leq\frac SE+\frac{\Delta^2}{2\pi\delta^2}\quad(S\geq0).
\]
\end{proposition}
\begin{proof}
Use $\TV(P,Q)=\sup_A|P(A)-Q(A)|$. For $m\geq0$,
\[
 \TV(\Normal(m,v),\Normal(0,v))
 =2\Phi(m/(2\sqrt v))-1\leq m/\sqrt{2\pi v}.
\]
Convexity under mixtures and $\E K=E$ bound the average-pair TV by
$\E[K\Delta/E]/\sqrt{2\pi v}=\Delta/\sqrt{2\pi v}$.
Both directed profiles at $\varepsilon\geq0$ are at most TV. Thus
$v=\Delta^2/(2\pi\delta^2)$ is feasible; adding the sampling term proves
the second inequality.
\end{proof}
The uniform bound is a feasibility estimate, not an asymptotic formula.
It shows why the transcript divergence cannot be attributed solely to the
number of releases: preserving the individual observations matters.

\section{Scope of the conclusions}
The finite theorem supplies a sufficient domain, not a complete classification
of privacy targets. It compares each smaller batch with full batch and does
not order every pair of smaller batches. The sparse cutoff is uniform in
energy, but fixes $E,\Delta,\varepsilon,\delta$. At a fixed dataset size $N$,
the expected batch $EN/T$ eventually falls below one; a joint limit in $N$
or a changing privacy target requires a separate analysis.

The fixed-rate asymptotic comparison in R\"ais\"a, J\"alk\"o, and
Honkela~\cite[Corollary 5.4]{raisa2024} follows a different path: for fixed
$q>0$, its calibrated raw-noise ratio satisfies
$\sigma(q,T)/(q\sigma(1,T))\to1$. It cannot be specialized to $q=E/T$
without uniformity in the rate. Here the expected participation, rather than
the sampling rate, is fixed.

All comparisons concern independent Poisson sampling, homogeneous independent
Gaussian noise, and equal-weight averages of fixed clipped contributions.
The common-channel argument requires the query to stay fixed across neighboring
datasets. It does not justify substituting a privately chosen optimization
trajectory into this model.

\clearpage
\appendix
\section{Calibration and the canonical experiment}\label{app:calibration}
\begin{proof}[Proof of Lemma~\ref{lem:calibration}]
We establish the finite-noise properties for arbitrary $T\geq1$ and
$0<q\leq1$, then specialize to full batch.

\paragraph{Monotonicity and feasibility.}
If $0\leq\sigma_1\leq\sigma_2$, adding independent centered Gaussian noise of
variance $\sigma_2^2-\sigma_1^2$ to each coordinate maps both transcript laws
at $\sigma_1$ to their counterparts at $\sigma_2$. Data processing for
hockey-stick divergence therefore makes the two-sided profile nonincreasing
in $\sigma$.

The sampled pair is also a common randomized postprocessing of the ordinary
Gaussian pair $\Normal(\Delta,\sigma^2)$ versus $\Normal(0,\sigma^2)$: retain
the input with probability $q$, and otherwise replace it with a fresh
$\Normal(0,\sigma^2)$ draw. Applying this channel independently gives
\[
 \AO_S(\sigma,\Delta,q,T,\varepsilon)
 \leq
 \AO_S(\sigma,\Delta,1,T,\varepsilon).
\]
The Gaussian product profile tends to zero as $\sigma\to\infty$, so the
feasible set is nonempty for every $\delta>0$.

\paragraph{Continuity and attainment.}
For $\sigma>0$, standardize the reference law to $T$ independent standard
normals and write
\[
 R_\sigma(z)=\prod_{i=1}^T
 \left[1-q+q\exp\left(\frac{\Delta z_i}{\sigma}
                         -\frac{\Delta^2}{2\sigma^2}\right)\right].
\]
With $\alpha=e^\varepsilon$, the directed profiles are
\begin{align*}
 H_\alpha(P_{q,\sigma}^{\otimes T}, Q_\sigma^{\otimes T})
 &=1-\E\min(R_\sigma,\alpha),\\
 H_\alpha(Q_\sigma^{\otimes T}, P_{q,\sigma}^{\otimes T})
 &=1-\alpha\E\min(R_\sigma,\alpha^{-1}).
\end{align*}
For fixed $z$, $R_\sigma(z)$ is continuous at positive $\sigma$. The truncated
integrands are bounded independently of $\sigma$, so bounded convergence
proves continuity there.

As $\sigma\downarrow0$, $R_\sigma(z)\to w:=(1-q)^T$ for every fixed $z$.
The same argument gives limits $1-w$ and $(1-\alpha w)_+$. At zero noise,
$Q_0^{\otimes T}$ is concentrated at the origin, and $P_{q,0}^{\otimes T}$
assigns it mass $w$. Direct optimization over events gives precisely these
two profiles. Their maximum is therefore continuous at zero as well.

The feasible set is closed and nonempty. A sequence of feasible noises
converging to its nonnegative infimum shows that the infimum is feasible.
If a positive minimum had profile strictly below $\delta$, continuity would
make a slightly smaller noise feasible. Hence the profile equals $\delta$
at every positive minimizing noise.

\paragraph{Full-batch calibration.}
At privacy variance $v>0$, each of the $E$ full-batch releases has variance
$Ev$. The product likelihood ratio is
\[
 \frac{dP^{\otimes E}}{dQ^{\otimes E}}(x)
 =\exp\left(\frac{\Delta}{Ev}\sum_{t=1}^E x_t
                     -\frac{\Delta^2}{2v}\right).
\]
It depends only on the sum, so the sum preserves both directed profiles.
After division by $E$, the two laws are $\Normal(\Delta,v)$ and
$\Normal(0,v)$. Their directed profiles agree by reflection, and the Gaussian
likelihood-threshold formula gives~\eqref{eq:gaussian-calibration}.

The identity
\[
 e^\varepsilon\phi(-a/2-\varepsilon/a)=\phi(a/2-\varepsilon/a)
\]
gives $G_\varepsilon'(a)=\phi(a/2-\varepsilon/a)>0$. Moreover,
$G_\varepsilon(a)\to0$ as $a\downarrow0$ and
$G_\varepsilon(a)\to1$ as $a\to\infty$, including when $\varepsilon=0$.
Thus the profile is continuous and strictly decreasing in $v>0$, from one
to zero. There is a unique positive $v_{\mathrm{FB}}$ with profile $\delta$,
and it is independent of $E$. Equivalently,
\begin{equation}\label{eq:fullbatch}
 \sigma_{E,E}=\kappa_{\varepsilon,\delta,\Delta}\sqrt E,
 \qquad V_E(E,S)=\kappa_{\varepsilon,\delta,\Delta}^2,
\end{equation}
where $\kappa_{\varepsilon,\delta,\Delta}$ is the one-release Gaussian
calibrated standard deviation. This is the analytical Gaussian calibration
of Balle and Wang~\cite{ballewang2018}.

\paragraph{Average-only calibration.}
For fixed $E,T$, write $p_k=\Pr\{K=k\}$ and $m_k=k\Delta/E$ in
\eqref{eq:average-experiment}. Adding independent Gaussian noise makes the
average-only profile nonincreasing in $v$, and
\eqref{eq:average-transcript-order} gives finite feasibility. At $v>0$,
standardizing the reference observation to $Z\sim\Normal(0,1)$ gives the
likelihood ratio
\[
 \overline R_v(z)=\sum_{k=0}^T p_k
       \exp\!\left(\frac{m_kz}{\sqrt v}-\frac{m_k^2}{2v}\right).
\]
This finite sum is pointwise continuous for $v>0$. As $v\downarrow0$, it
converges pointwise to $p_0=(1-E/T)^T$: every term with $k>0$ tends to zero.
The same bounded-truncation formulas used above prove continuity of both
directed profiles, including at zero. At zero the source has mass $p_0$
at the reference atom and the remaining mass at positive locations, giving
the same directed endpoint values $1-p_0$ and $(1-e^\varepsilon p_0)_+$.
Closedness and the same minimum argument prove attainment and equality to
the target at a positive minimizer. At $T=E$, $K=E$ almost surely, so the
average pair is $\Normal(\Delta,v)$ versus $\Normal(0,v)$ and
$\overline v_{E,E}=v_{\mathrm{FB}}$.
\end{proof}

\subsection{Domination for fixed recordwise queries}\label{app:canonical-domination}
For completeness, consider a fixed recordwise query whose contributions
have magnitude at most $\Delta$. Under add/remove adjacency, let $d$ be the
differing contribution, so $|d|\leq\Delta$. The Gaussian channel
\[
 x\longmapsto\frac d\Delta x+\zeta,
 \qquad
 \zeta\sim\Normal\left(0,\sigma^2\left(1-\frac{d^2}{\Delta^2}\right)\right)
\]
maps $\Normal(0,\sigma^2)$ to itself and $\Normal(\Delta,\sigma^2)$ to
$\Normal(d,\sigma^2)$. It therefore maps the canonical one-release pair to
the pair with differing contribution $d$. The statement also holds at
$\sigma=0$, with degenerate Gaussian noise.

Apply this channel independently across releases, then add the vector of
shared sampled contributions. That vector is independent of the differing
record's inclusion indicators and of the Gaussian noise, and it has the
same law on the two neighboring datasets. These common channels produce
the clipped-sum transcript, so data processing bounds its profile by the
canonical profile. The query must be held fixed across the neighboring
comparison; a dataset-dependent choice of the query would require a
separate argument.

For the average-only comparison, let
$C_T=E^{-1}\sum_{t=1}^T\sum_{i\,\mathrm{shared}} B_{it}\widetilde g_i$
be the shared sampled contribution to the average. It is independent of
the differing record's participation count and of the added noise. Apply
directly to the canonical average the common channel
\[
 X\longmapsto\frac d\Delta X+\zeta+C_T,\qquad
 \zeta\sim\Normal\!\left(0,v\left(1-\frac{d^2}{\Delta^2}\right)\right),
\]
where $\zeta$ and $C_T$ are independent of $X$ and of each other. Under the
reference law the output is $C_T+\Normal(0,v)$. Conditional on $K=k$ under
the source law, it is $C_T+kd/E+\Normal(0,v)$: scaling the input noise and
adding $\zeta$ leaves Gaussian variance $v$. These are exactly the two
neighboring averages when each raw release receives Gaussian noise of
variance $E^2v/T$. The construction
also holds at $v=0$. Data processing therefore bounds each such neighboring
average profile by $\overline{\mathcal A}_{E,T}(v;\varepsilon)$.

\section{Binomial sign comparisons}\label{app:binomial-signs}
Throughout this appendix, $T>E>0$ are integers,
$p=E/T$, $r=1-p$, $K\sim\operatorname{Binomial}(T,p)$, and $Z=K-E$.
The gap $\Gamma_{E,T}$ is defined in~\eqref{eq:main-count-gap}. Its relevant
moments are
\begin{equation}\label{eq:binomial-moments}
 \E Z=0,\qquad \E Z^2=Er,\qquad \E Z^3=Er(1-2p).
\end{equation}
These follow by expanding the sum of $T$ independent centered Bernoulli
variables: every mixed term containing a singleton centered factor vanishes.
Both sign arguments below also use the size-bias identities: if
$K\sim\operatorname{Binomial}(T,p)$ and $L\sim\operatorname{Binomial}(T-1,p)$,
then for every integer $k$,
\begin{equation}\label{eq:binomial-size-bias}
 k\Pr\{K=k\}=Tp\Pr\{L=k-1\},\qquad
 (T-k)\Pr\{K=k\}=Tr\Pr\{L=k\}.
\end{equation}
Adding them gives
$\Pr\{K=k\}=r\Pr\{L=k\}+p\Pr\{L=k-1\}$.

\begin{theorem}[The binomial center comparison]\label{thm:rate-adaptive-center}
Let $T>E>0$ be integers and $h>0$. Then
$\Gamma_{E,T}(\theta,h)>0$ under the following conditions:
\begin{equation}\label{eq:rate-adaptive-center-condition}
 \begin{array}{c|c}
 T>2E & \theta\geq c_{E,T}:=(1-2E/T)/3,\\
 T=2E & \theta>0,\\
 E<T<2E & \theta\geq0.
 \end{array}
\end{equation}
At $T=2E$ and $\theta=0$, the gap is zero.
\end{theorem}

We first show that a nonnegative gap becomes positive when the center is
increased. We then establish the starting comparison at $c_{E,T}$ in the
low-rate case and at zero in the high-rate case. Symmetry handles the
half-rate case.

\subsection{Propagation in the center}
\begin{lemma}[Propagation of positivity]\label{lem:binomial-propagation}
Fix $h>0$ and $0\leq\theta_1<\theta_2$. If
$\Gamma_{E,T}(\theta_1,h)\geq0$, then
$\Gamma_{E,T}(\theta_2,h)>0$.
\end{lemma}

\begin{proof}
For $t\in\R$, put $k_\theta(t)=h\phi(h(t-\theta))$. For every integer $z$,
\[
 \Phi(h(z-\theta))-\Phi(-h\theta)=\int_0^z k_\theta(t)\,dt.
\]
Let $A_\theta$ be the expectation of the positive integrals and $B_\theta$
the expectation of the absolute values of the negative integrals. Then
$\Gamma_{E,T}(\theta,h)=A_\theta-B_\theta$, and
\[
 \frac{k_{\theta_2}(t)}{k_{\theta_1}(t)}
 =\chi\exp\bigl(h^2(\theta_2-\theta_1)t\bigr),
 \qquad
 \chi=\exp\left(-\frac{h^2}{2}(\theta_2^2-\theta_1^2)\right)>0.
\]
The ratio is strictly larger than $\chi$ for $t>0$ and strictly smaller for
$t<0$. The centered nondegenerate binomial law has positive mass on both
sides of zero, so
\[
 A_{\theta_2}>\chi A_{\theta_1},\qquad
 B_{\theta_2}<\chi B_{\theta_1}.
\]
Subtracting yields
$\Gamma_{E,T}(\theta_2,h)>\chi\Gamma_{E,T}(\theta_1,h)\geq0$.
\end{proof}

\subsection{Low rate: the critical center}
Assume $T>2E$ and abbreviate $c=c_{E,T}=(1-2p)/3$, so $0<c<1/3$.
By~\eqref{eq:gap-small-scale}, the cubic term of $\Gamma_{E,T}(c,h)$
vanishes at this center, so the cubic expansion does not decide the sign.
We need an argument valid at every scale. Since $\Gamma_{E,T}(c,0)=0$, it is
enough to show that $\Gamma_{E,T}(c,\cdot)$ is strictly increasing on
$(0,\infty)$. Differentiation over the finite support gives
\begin{equation}\label{eq:binomial-scale-derivative}
 \frac{d}{dh}\Gamma_{E,T}(c,h)=\phi(0)D(h^2/2),
\end{equation}
where
\begin{equation}\label{eq:binomial-laplace}
 D(s)=\E[(Z-c)e^{-s(Z-c)^2}]+ce^{-sc^2}.
\end{equation}

Each support value enters $D$ through a signed weight $z-c$ and a squared
distance $(z-c)^2$. Accordingly, place an atom of weight
$\nu_z=\Pr\{Z=z\}(z-c)$ at $y_z=(z-c)^2$ for each support value $z$, and one
more atom of weight $c$ at $c^2$ for the second term. Then
$D(s)=\sum\nu e^{-sy}$ is the Laplace transform of the finite signed
measure~$\nu$. Its first two moments are the coefficients that vanished in the
expansion:
\begin{equation}\label{eq:binomial-cancellations}
 \sum\nu=\E Z=0,\qquad
 \sum\nu y=\E[(Z-c)^3]+c^3=\E Z^3-3c\E Z^2=0,
\end{equation}
so $D(0)=0$ because the count has mean $E$, and $D'(0)=0$ because of the
choice of $c$.

The positive and negative parts of $\nu$ therefore have equal mass and equal
first moment. $D(s)$ compares their integrals of the convex function
$y\mapsto e^{-sy}$, so the sign should come from which part is more spread
out. For measures of equal mass and first moment, spread is measured by the
hinge transform
\begin{equation}\label{eq:binomial-double-cumulative}
 J(t)=\sum\nu(t-y)_+,\qquad (u)_+=\max\{u,0\}:
\end{equation}
$J\geq0$ says that the positive part dominates the negative part in convex
order. We use $J$ only through an exact identity, which also yields
strictness. For $s>0$, finite summation and
\[
 e^{-sy}=s^2\int_y^\infty e^{-st}(t-y)\,dt
\]
give
\begin{equation}\label{eq:binomial-laplace-hinge}
 D(s)=s^2\int_0^\infty e^{-st}J(t)\,dt.
\end{equation}
By~\eqref{eq:binomial-cancellations}, $J(t)=t\sum\nu-\sum\nu y=0$ beyond
the largest location, so $J$ is continuous, piecewise affine and compactly
supported. The low-rate comparison thus reduces to showing that $J\geq0$
and $J\not\equiv0$.

To locate the kinks of $J$, note that the atoms at $c^2$ from $Z=0$ and from
the extra term combine to net weight $c(1-\Pr\{Z=0\})>0$. The other positive
weights come from displacements $z=j\geq1$ and sit at the \emph{positive
knots} $(j-c)^2$. The negative weights come from $z=-j\leq-1$ and sit at
$(j+c)^2$. Because $0<c<1/2$, the locations interlace:
\[
 c^2<(1-c)^2<(1+c)^2<(2-c)^2<(2+c)^2<\cdots.
\]
Between consecutive positive knots, $J$ has at most one kink, and its slope
there drops. So $J$ is concave on each such interval, and it suffices to check
$J\geq0$ at the positive knots. At a positive knot $t=(j-c)^2$, $J(t)$ is a
truncated cubic expectation over $|z-c|<j-c$. For $j\leq E$ we sum it
exactly. Put $C=T-E$, let $L\sim\operatorname{Binomial}(T-1,p)$, and write
$\ell_k=\Pr\{L=k\}$, with zero mass outside its support.
By~\eqref{eq:binomial-size-bias},
\[
 r(E+z)\Pr\{Z=z\}=Er\,\ell_{E+z-1},\qquad
 p(C-z)\Pr\{Z=z\}=Er\,\ell_{E+z}.
\]
Hence any summand of the form $r(E+z)A(z)-p(C-z)A(z+1)$, weighted by
$\Pr\{Z=z\}$, becomes $Er\{\ell_{E+z-1}A(z)-\ell_{E+z}A(z+1)\}$. Summed over
consecutive integers, it telescopes to two boundary terms. The expression
$r(E+z)A(z)-p(C-z)A(z+1)$ is the binomial Stein operator applied to $A$.
Our summand is cubic in $z$, so we look for a quadratic $A$.

\begin{lemma}[A binomial telescoping identity]\label{lem:binomial-stein-polynomial}
Put $C=T-E$. For any real $R$, define
\[
 A_R(z)=R^2-2Er-3c^2+z-z^2.
\]
Then for every $z\in\R$,
\begin{equation}\label{eq:binomial-cubic-stein-polynomial}
 \begin{split}
 &(z-c)\{R^2-(z-c)^2\}+c(R^2-c^2)\\
 &\hspace{2em}=r(E+z)A_R(z)-p(C-z)A_R(z+1).
 \end{split}
\end{equation}
\end{lemma}
\begin{proof}
The size-bias identities suggest the operator on the right: after weighting
by $\Pr\{Z=z\}$ it becomes a difference of consecutive predecessor-binomial
terms. To match the cubic on the left, try
$A(z)=\alpha+\beta z-z^2$. Using $pC=Er$ and $p+r=1$, its coefficients are
\[
 r(E+z)A(z)-p(C-z)A(z+1)
 =-z^3+(\beta-2p)z^2
   +(\alpha+p\beta-p+2Er)z+Er(1-\beta).
\]
The left side of~\eqref{eq:binomial-cubic-stein-polynomial} expands to
$-z^3+3cz^2+(R^2-3c^2)z$.
Matching the quadratic term gives $\beta=2p+3c=1$, because
$c=(1-2p)/3$. The constant term then vanishes, and matching the linear term
gives $\alpha=R^2-2Er-3c^2$. This proves the identity for arbitrary $R$.
\end{proof}

\begin{lemma}[Nonnegativity of the hinge transform]
\label{lem:binomial-double-cumulative}
If $T>2E$, then $J(t)\geq0$ for every $t\geq0$, and
$J((1-c)^2)>0$.
\end{lemma}

\begin{proof}
\emph{Positive knots with $j\leq E$.} Fix $1\leq j\leq E$ and
$t=R_j^2=(j-c)^2$, and abbreviate $A_j=A_{R_j}$.
The atom at $y_z$ is active, $(t-y_z)_+>0$, exactly when
$|z-c|<j-c$, that is, $2c-j<z<j$; since $0<2c<1$, this means
$1-j\leq z\leq j-1$. The extra atom at $c^2$ is active as well. Splitting its
weight as $c=c\Pr\{|Z|\leq j-1\}+c\Pr\{|Z|\geq j\}$ puts the first part inside
the sum:
\[
 J((j-c)^2)=\sum_{z=1-j}^{j-1}\Pr\{Z=z\}
 \bigl[(z-c)\{R_j^2-(z-c)^2\}+c(R_j^2-c^2)\bigr]
 +c(R_j^2-c^2)\Pr\{|Z|\geq j\}.
\]
The bracket is the left side of~\eqref{eq:binomial-cubic-stein-polynomial}.
By Lemma~\ref{lem:binomial-stein-polynomial} and the size-bias identities,
the sum telescopes:
\begin{align*}
 &Er\sum_{z=1-j}^{j-1}
 \{\ell_{E+z-1}A_j(z)-\ell_{E+z}A_j(z+1)\}\\
 &\hspace{2em}=Er\{\ell_{E-j}A_j(1-j)-\ell_{E+j-1}A_j(j)\}.
\end{align*}
Both endpoint values equal
\[
 A_j(1-j)=A_j(j)=b_j:=j(1-2c)-2Er-2c^2,
\]
and $R_j^2-c^2=j(j-2c)$. Thus
\begin{equation}\label{eq:binomial-positive-knot}
 \begin{split}
 J((j-c)^2)={}&Er b_j(\ell_{E-j}-\ell_{E+j-1})\\
              &+cj(j-2c)\Pr\{|Z|\geq j\}.
 \end{split}
\end{equation}
The tail term is nonnegative. For the boundary term, $1-2c>0$ gives
$b_j\leq b_E$, and the identity $1-2r=-3c$ gives
\[
 b_E=-5Ec-2c^2<0.
\]
So the boundary term is nonnegative as soon as the lower mass of each pair
is the smaller one:
\begin{equation}\label{eq:binomial-pairing}
 \ell_{E-j}\leq\ell_{E+j-1}\qquad(1\leq j\leq E).
\end{equation}
At $j=1$ the two masses are equal. For $1\leq j<E$, the ratios taking the
upper and lower members of one pair to the next are
\[
 \frac{\ell_{E+j}}{\ell_{E+j-1}}
 =\frac{C-j}{E+j}\frac EC,\qquad
 \frac{\ell_{E-j-1}}{\ell_{E-j}}
 =\frac{E-j}{C+j}\frac CE.
\]
The first ratio is at least the second: after clearing positive
denominators, their difference has the sign of $j^2(C^2-E^2)\geq0$, since
$C>E$. Induction proves~\eqref{eq:binomial-pairing}, so
$J((j-c)^2)\geq0$ for $1\leq j\leq E$. At $j=1$ the boundary term vanishes,
while the tail term is $c(1-2c)\Pr\{|Z|\geq1\}>0$. Hence $J((1-c)^2)>0$.

\emph{Positive knots with $j>E$.} Now every negative displacement satisfies
$|z-c|\leq E+c<j-c$, so all atoms except the positive ones with $z\geq j$ are
active. Subtracting the vanishing full sum
$\sum\nu(t-y)=t\sum\nu-\sum\nu y=0$ leaves
\[
 J((j-c)^2)=\sum_{z\geq j}\Pr\{Z=z\}(z-c)
             \bigl((z-c)^2-(j-c)^2\bigr)\geq0.
\]

\emph{Between the knots.} Before $c^2$, $J$ vanishes. On
$[c^2,(1-c)^2]$ only the net atom at $c^2$ is active, so
\[
 J(t)=c(1-\Pr\{Z=0\})(t-c^2)\geq0
 \qquad(c^2\leq t\leq(1-c)^2).
\]
Between $(j-c)^2$ and $(j+1-c)^2$, $j\geq1$, the only possible interior knot
is $(j+c)^2$, with weight $(-j-c)\Pr\{Z=-j\}\leq0$. So $J$ is concave there
and lies above the chord joining its nonnegative endpoint values. Beyond the
largest location $J$ vanishes. This proves $J(t)\geq0$ for all $t\geq0$.
\end{proof}

By Lemma~\ref{lem:binomial-double-cumulative}, $J$ is continuous,
nonnegative, and positive at $(1-c)^2$. Hence the integral
in~\eqref{eq:binomial-laplace-hinge} is strictly positive for every $s>0$.
Equation~\eqref{eq:binomial-scale-derivative} and
$\Gamma_{E,T}(c,0)=0$ now give
\begin{equation}\label{eq:binomial-low-boundary-strict}
 \Gamma_{E,T}(c,h)>0\qquad(T>2E,\ h>0).
\end{equation}

\subsection{Zero-center signs and reflection}
\begin{lemma}[The zero-center sign]\label{lem:binomial-reflected-zero}
For $h>0$, the gap $\Gamma_{E,T}(0,h)$ is negative when $T>2E$, zero when
$T=2E$, and positive when $E<T<2E$.
\end{lemma}

\begin{proof}
We first treat a general low-rate mean $m$, so that the same computation also
serves the high-rate case after reflection. Fix an integer $m$ with
$0<m<T/2$. Let $K\sim\operatorname{Binomial}(T,m/T)$ and
$L\sim\operatorname{Binomial}(T-1,m/T)$, and put
\[
 \psi_h(k)=\exp\left(-\frac{h^2(k-m)^2}{2}\right).
\]
The size-bias identities~\eqref{eq:binomial-size-bias} with $p=m/T$ give,
for every function $f$,
\[
 \E[(K-m)f(K)]
 =m(1-m/T)\E[f(L+1)-f(L)].
\]
Differentiating the gap and applying this identity gives
\begin{equation}\label{eq:binomial-reflected-stein}
 \frac{d}{dh}\Gamma_{m,T}(0,h)
 =\phi(0)m(1-m/T)\E[\psi_h(L+1)-\psi_h(L)].
\end{equation}
Pair $L=m-1-j$ with $L=m+j$ for $0\leq j<m$. Their increments are $d_j$
and $-d_j$, respectively, where
\[
 d_j=e^{-h^2j^2/2}-e^{-h^2(j+1)^2/2}>0.
\]
By~\eqref{eq:binomial-pairing}, applied with $E=m$, the upper mass is at
least the lower mass, so each pair contributes a nonpositive amount.
The remaining upper atoms, indexed by $m\leq j<T-m$, have positive mass
and negative increment $-d_j$. At least one exists because $m<T/2$.
Thus the derivative in~\eqref{eq:binomial-reflected-stein} is negative for
$h>0$. Since $\Gamma_{m,T}(0,0)=0$, the low-rate gap is negative.

If $T=2E$, the law of $Z$ is symmetric and Gaussian symmetry gives
$\Gamma_{E,T}(0,h)=0$. If $E<T<2E$, put $m=T-E<T/2$ and reflect the count:
$K'=T-K\sim\operatorname{Binomial}(T,m/T)$, with $K'-m=-Z$. Hence
\[
 \Gamma_{E,T}(0,h)=-\Gamma_{m,T}(0,h)>0.
\]
\end{proof}

\begin{proof}[Proof of Theorem~\ref{thm:rate-adaptive-center}]
For $T>2E$, equation~\eqref{eq:binomial-low-boundary-strict} proves positivity
at $\theta=c_{E,T}$, and Lemma~\ref{lem:binomial-propagation} extends it to
larger centers. For $E<T<2E$, Lemma~\ref{lem:binomial-reflected-zero} gives
positivity at zero, and the same propagation lemma handles every positive
center. At $T=2E$, symmetry gives equality at zero and propagation gives
strict positivity for $\theta>0$.
\end{proof}

\subsection{Strict comparison at the symmetric endpoint}
At the symmetric endpoint of Theorem~\ref{thm:rate-adaptive-pointwise}, the
full-batch threshold gives a tie. A slightly larger threshold already
gives a strict improvement using the average alone.

\begin{lemma}[Improving the threshold at half rate]\label{lem:strict-average-endpoint}
Let $E\in\N_{>0}$, $a>0$, and $K\sim\operatorname{Binomial}(2E,1/2)$.
Let $Q=\Normal(0,1)$, and let $P$ be the mixture whose conditional law
given $K=k$ is $\Normal(ak/E,1)$. With $\varepsilon=a^2/2$,
\[
 H_{e^\varepsilon}(P, Q)
 >\frac12-e^\varepsilon\overline\Phi(a).
\]
\end{lemma}

\begin{proof}
Symmetry of $K-E$ and of the Gaussian noise makes $P$ symmetric about
$a$, with no atom there. Thus $P((a,\infty))=1/2$, and the value of this
ray is $\delta_*:=1/2-e^\varepsilon\overline\Phi(a)$.
The likelihood ratio of the mixture is the continuous function
\[
 L(y)=\E\exp\!\left(\frac{aK}{E}y-\frac{a^2K^2}{2E^2}\right).
\]
Completing the square at the threshold gives
\[
 L(a)=e^{a^2/2}\E\exp\!\left[
       -\frac{a^2}{2}(K/E-1)^2\right]<e^{a^2/2}=e^\varepsilon.
\]
The inequality is strict because $a>0$ and $\Pr\{K\ne E\}>0$.
By continuity there is $\eta>0$ such that $L(y)<e^\varepsilon$ for
$a\leq y\leq a+\eta$. Since $\phi(y)>0$,
\[
 P((a+\eta,\infty))-e^\varepsilon Q((a+\eta,\infty))
 =\delta_*-\int_a^{a+\eta}(L(y)-e^\varepsilon)\phi(y)\,dy
 >\delta_*.
\]
Taking the supremum over events proves the claim.
\end{proof}

\subsection{Sharpness uniformly over the Gaussian scale}
Equation~\eqref{eq:gap-small-scale} identifies the cubic coefficient.
Since $\operatorname{Var}(Z)>0$, every $\theta<c_{E,T}$ gives a negative gap
for sufficiently small $h>0$. Together with
Theorem~\ref{thm:rate-adaptive-center}, this proves that, in the low-rate
regime, $\theta\geq c_{E,T}$ is necessary and sufficient for positivity at
every $h>0$. This is sharpness of the scale-uniform sum-event comparison,
not a classification of the full transcript profile when the scale and
privacy parameters are coupled.

\section{The fixed-participation privacy limit}\label{app:sparse-limit}
We first prove the profile limit for $Tq_T\to\rho>0$, then invert it to
obtain the calibrated-noise offset.

\begin{theorem}[The fixed-participation profile limit]\label{thm:accountant-limit}
Let $\rho>0$, $\Delta>0$, and $\varepsilon\geq0$. Suppose
$q_T\in[0,1]$, $Tq_T\to\rho$, $\sigma_T>0$ eventually, and
\[
 \frac{\Delta}{\sigma_T}-\sqrt{2\log T}\longrightarrow s\in\R.
\]
Then
\[
 \AO_S(\sigma_T,\Delta,q_T,T,\varepsilon)
 \longrightarrow1-\exp\{-\rho\Phi(s)\}.
\]
\end{theorem}

\begin{proof}
Put $\mu_T=\Delta/\sigma_T$ and $b_T=\sqrt{2\log T}$. Standardizing every
coordinate by $\sigma_T$ makes the reference observations independent
$Z_i\sim\Normal(0,1)$. Their likelihood ratio is
\begin{equation}\label{eq:sparse-likelihood-ratio}
 R_T=\prod_{i=1}^T
 \left[1-q_T+q_T\exp\left(\mu_TZ_i-\frac{\mu_T^2}{2}\right)\right].
\end{equation}

Write $Y_T(z)=\exp(\mu_Tz-\mu_T^2/2)$. We linearize $\log R_T$ after
truncating each coordinate at a level $t_T$.
The level has two jobs. It must be high enough that, with reference
probability tending to one, no coordinate exceeds it, so truncation changes
nothing. We therefore need $T\overline\Phi(t_T)\to0$; Mills' bound makes
$T/t_T=o(e^{t_T^2/2})$ sufficient. The level must also be low enough that every
perturbation $q_T(Y_T-1)$ of the likelihood ratio is uniformly small and the
perturbations have vanishing total second moment. Below this level,
$Y_T\leq e^{t_T^2/2}$; together with the moment bound below, this makes
$e^{t_T^2/2}=o(T)$ sufficient. The naive level $t_T=b_T$ gives
$e^{t_T^2/2}=T$, and our coarse second-moment bound stays of order one.
Lowering the level by $\log b_T/(2b_T)$ gives
$e^{t_T^2/2}=Tb_T^{-1/2}(1+o(1))$, which lies between $T/t_T$ and $T$. The first
truncated moment then produces the limit, and the second moment controls both
the fluctuation of the linear term and the linearization error.

\paragraph{The reference maximum.}
For all sufficiently large $T$, define
\[
 t_T=b_T-\frac{\log b_T}{2b_T}.
\]
Mills' upper bound gives
\[
 T\Pr\{Z>t_T\}\leq T\frac{\phi(t_T)}{t_T}\longrightarrow0,
\]
because
\[
 \log T-\frac{t_T^2}{2}
 =\frac{\log b_T}{2}-\frac{(\log b_T)^2}{8b_T^2}.
\]
The upper bound is of order $b_T^{-1/2}$. Thus every coordinate lies below
$t_T$ with probability tending to one. We will also use
\begin{equation}\label{eq:truncation-scales}
 \frac{e^{t_T^2/2}}{T}=b_T^{-1/2}(1+o(1)),\qquad
 q_Te^{t_T^2/2}\to0,\qquad Tq_T^2e^{t_T^2/2}\to0.
\end{equation}
The last two estimates follow by multiplying the first by $Tq_T\to\rho$
and by $(Tq_T)^2$, respectively.

\paragraph{Truncated moments.}
Define
\[
 X_T(z)=q_T(Y_T(z)-1)\mathbf1_{\{z\leq t_T\}}.
\]
Completing the square gives
\begin{align}
 \E[Y_T(Z)\mathbf1_{\{Z\leq t_T\}}]
 &=\Phi(t_T-\mu_T),\label{eq:tilted-first}\\
 \E[Y_T(Z)^2\mathbf1_{\{Z\leq t_T\}}]
 &=e^{\mu_T^2}\Phi(t_T-2\mu_T).\label{eq:tilted-second}
\end{align}
Hence
\[
 T\E X_T(Z)=Tq_T[\Phi(t_T-\mu_T)-\Phi(t_T)]
 \longrightarrow-\rho\Phi(s),
\]
since $t_T-\mu_T\to-s$ and $\Phi(t_T)\to1$.

Because $Y_T\geq0$, $(Y_T-1)^2\leq Y_T^2+1$, so
\begin{equation}\label{eq:truncated-second-bound}
 T\E X_T(Z)^2\leq Tq_T^2
 \left[e^{\mu_T^2}\Phi(t_T-2\mu_T)+\Phi(t_T)\right].
\end{equation}
The second term tends to zero. For the first,
$2\mu_T-t_T=b_T+2s+o(1)\to\infty$, and eventually $2\mu_T-t_T\geq1$.
Mills' bound and the identity
\[
 \mu_T^2-\frac{(2\mu_T-t_T)^2}{2}
 =\frac{t_T^2}{2}-(\mu_T-t_T)^2
\]
therefore imply
\[
 e^{\mu_T^2}\Phi(t_T-2\mu_T)
 \leq\frac{e^{t_T^2/2}}{\sqrt{2\pi}}.
\]
By~\eqref{eq:truncation-scales}, the right side
of~\eqref{eq:truncated-second-bound} tends to zero. Independence gives
\[
 \operatorname{Var}\left(\sum_{i=1}^T X_T(Z_i)\right)
 \leq T\E X_T(Z)^2\longrightarrow0.
\]
Chebyshev's inequality yields
\begin{equation}\label{eq:linear-sum-limit}
 \sum_{i=1}^T X_T(Z_i)\longrightarrow-\rho\Phi(s)
 \quad\text{in reference probability}.
\end{equation}

\paragraph{The log likelihood ratio.}
Let $A_T=\{\max_i Z_i\leq t_T\}$, so $\Pr(A_T^c)\to0$.
Eventually $\mu_T\geq0$, and for $z\leq t_T$,
\[
 |q_T(Y_T(z)-1)|
 \leq q_T+q_T\exp\left(\mu_Tt_T-\frac{\mu_T^2}{2}\right)
 \leq q_T+q_Te^{t_T^2/2}\longrightarrow0.
\]
The bound is uniform in $z\leq t_T$. Eventually every truncated perturbation
therefore has absolute value at most $1/2$, where
$|\log(1+x)-x|\leq x^2$. On $A_T$,
\begin{align*}
 \mathbf1_{A_T}\left|\log R_T-\sum_{i=1}^T X_T(Z_i)\right|
 &\leq\mathbf1_{A_T}\sum_{i=1}^T
   |\log(1+X_T(Z_i))-X_T(Z_i)|\\
 &\leq\sum_{i=1}^T X_T(Z_i)^2.
\end{align*}
The expectation of the last expression tends to zero by
\eqref{eq:truncated-second-bound}. Markov's inequality, the bound on
$\Pr(A_T^c)$, and~\eqref{eq:linear-sum-limit} imply
\begin{equation}\label{eq:log-likelihood-limit}
 \log R_T\longrightarrow-\rho\Phi(s)
 \quad\text{in reference probability}.
\end{equation}
Thus $R_T\to c_s:=e^{-\rho\Phi(s)}$ in reference probability.
Although $\E R_T=1$ for every $T$, $c_s<1$: rare large values retain the
missing mean, so convergence of unbounded expectations cannot be used.

\paragraph{The directed privacy limits.}
Let $\alpha=e^\varepsilon\geq1$. Writing $P_T,Q_T$ for the two transcript
laws, the directed profiles are
\begin{align*}
 H_\alpha(P_T, Q_T)&=1-\E_{Q_T}\min\{R_T,\alpha\},\\
 H_\alpha(Q_T, P_T)&=1-\alpha\E_{Q_T}\min\{R_T,\alpha^{-1}\}.
\end{align*}
Convergence in probability to a constant implies convergence of expectations
for bounded continuous functions. Applying this fact to the two truncations
gives
\[
 H_\alpha(P_T, Q_T)\to1-c_s,\qquad
 H_\alpha(Q_T, P_T)\to(1-\alpha c_s)_+.
\]
Since $(1-\alpha c_s)_+\leq1-c_s$, taking the maximum proves the theorem.
\end{proof}

\begin{theorem}[Fixed-participation calibration]
\label{thm:fixed-participation-calibration}
Let $\rho>0$, $\Delta>0$, $\varepsilon\geq0$, and
$0<\delta<1-e^{-\rho}$. Suppose $q_T\in[0,1]$ and $Tq_T\to\rho$, and define
\[
 \sigma_T^{\mathrm{cal}}:=\inf\left\{\sigma\geq0:
 \AO_S(\sigma,\Delta,q_T,T,\varepsilon)\leq\delta\right\}.
\]
Then $\sigma_T^{\mathrm{cal}}>0$ eventually, and
\[
 \frac{\Delta}{\sigma_T^{\mathrm{cal}}}-\sqrt{2\log T}
 \longrightarrow s_{\rho,\delta}:=
 \Phi^{-1}\!\left(-\frac{\log(1-\delta)}\rho\right).
\]
For the capped exact-budget schedule $q_T=\min\{1,\rho/T\}$, it follows that
\[
 \frac{\sigma_T^{\mathrm{cal}}\sqrt{2\log(1/q_T)}}\Delta\longrightarrow1.
\]
\end{theorem}

\begin{proof}
The function $L_\rho(s)=1-e^{-\rho\Phi(s)}$ is continuous and strictly
increasing, with range $(0,1-e^{-\rho})$ and inverse value
$L_\rho(s_{\rho,\delta})=\delta$. For a fixed offset $u$, consider the
eventually positive candidate
\[
 \sigma_T(u)=\frac\Delta{b_T+u},\qquad b_T=\sqrt{2\log T}.
\]
Given $\eta>0$,
\[
 L_\rho(s_{\rho,\delta}-\eta)<\delta
 <L_\rho(s_{\rho,\delta}+\eta).
\]
Theorem~\ref{thm:accountant-limit} makes the higher-noise candidate
$\sigma_T(s_{\rho,\delta}-\eta)$ eventually feasible and the lower-noise
candidate $\sigma_T(s_{\rho,\delta}+\eta)$ eventually infeasible.
The finite calibration properties proved in Appendix~\ref{app:calibration}
apply once $q_T>0$, which holds eventually. Monotonicity and feasibility give
\[
 \frac\Delta{b_T+s_{\rho,\delta}+\eta}
 \leq\sigma_T^{\mathrm{cal}}
 \leq\frac\Delta{b_T+s_{\rho,\delta}-\eta}.
\]
In particular, the calibrated noise is eventually positive. Taking
reciprocals gives
\[
 s_{\rho,\delta}-\eta
 \leq\frac\Delta{\sigma_T^{\mathrm{cal}}}-b_T
 \leq s_{\rho,\delta}+\eta.
\]
Since $\eta$ is arbitrary, the offset limit follows. For the capped schedule,
$q_T=\rho/T$ eventually and $\log(1/q_T)/\log T\to1$, giving the final claim.
\end{proof}

Equivalently, the calibrated standard deviation satisfies
\begin{equation}\label{eq:offset-approximation}
 \sigma_T^{\mathrm{cal}}
 =\frac\Delta{\sqrt{2\log T}+s_{\rho,\delta}+o(1)}.
\end{equation}

\begin{proof}[Proof of Theorem~\ref{thm:smallbatch}]
Apply Theorem~\ref{thm:fixed-participation-calibration} with $\rho=E$.
For $T\geq E$, its capped schedule is exactly $q=E/T$, so the offset limit
and~\eqref{eq:sigmaasymptotic} follow. Squaring and multiplying by $T/E^2$
gives~\eqref{eq:privacydivergence}. The sampling term in~\eqref{eq:variance}
converges to $S/E$, while the privacy variance diverges. Their sum is
therefore asymptotic to $v_{E,T}$, proving~\eqref{eq:total-equivalent}.
\end{proof}

\section{Rational bounds for the finite reversal}\label{app:rational-calculation}
This appendix proves the privacy bounds used in
Theorem~\ref{thm:strict-interior-counterexample}. We partition the first two
reference coordinates into rectangles and integrate over the third
coordinate exactly. Rational interval arithmetic then bounds both directed
profiles. The supplement~\cite{rational-supplement} contains the individual
logarithm brackets and unrounded interval expressions. Decimal values, where
given, are for orientation only.

\subsection{Conditional tails and rectangle bounds}
Fix $E=2$, $T=3$, $\Delta=1$, and $\varepsilon=1/4$. The candidate
privacy variance $v=5/32$ corresponds to
\[
 \sigma^2=\frac{E^2v}{T}=\frac5{24},\qquad
 \mu=\frac\Delta\sigma=\sqrt{\frac{24}{5}}.
\]
Write $P=P_{2/3,\sqrt{5/24}}$, $Q=Q_{\sqrt{5/24}}$, $q=2/3$, and
$\alpha=e^{1/4}$. After standardization, the one-coordinate likelihood ratio
under the reference law is
\[
 r(z)=\frac13+\frac23\exp\left(\mu z-\frac{\mu^2}{2}\right).
\]
Condition on the first two reference coordinates and put
$A=r(Z_1)r(Z_2)$. The remaining conditional hinges are
\[
 F(A)=\E[(Ar(Z)-\alpha)_+],\qquad
 G(A)=\E[(1-\alpha Ar(Z))_+],
\]
where $Z$ is standard normal. If $\alpha>A(1-q)$, define
\[
 c_F(A)=\frac{\log((\alpha/A-(1-q))/q)+\mu^2/2}{\mu}.
\]
Gaussian tail integration gives
\[
 F(A)=\bigl(A(1-q)-\alpha\bigr)\overline\Phi(c_F(A))
       +Aq\,\overline\Phi(c_F(A)-\mu).
\]
If $\alpha\leq A(1-q)$, then $F(A)=A-\alpha$. Similarly, if
$1>\alpha A(1-q)$, define
\[
 c_R(A)=\frac{\log((1/(\alpha A)-(1-q))/q)+\mu^2/2}{\mu}.
\]
Then
\[
 G(A)=\bigl(1-\alpha A(1-q)\bigr)\Phi(c_R(A))
       -\alpha Aq\,\Phi(c_R(A)-\mu),
\]
whereas $G(A)=0$ if $1\leq\alpha A(1-q)$.

Both $F$ and $G$ are convex in $A$, being expectations of positive parts of
affine functions. Let $\gamma=\Normal(0,1)$. On a rectangle $C$ where
$L\leq A\leq U$ with $L<U$, convexity gives
\begin{equation}\label{eq:rectangle-chord}
 \int_C F(A)\,d\gamma^{\otimes2}
 \leq\frac{UQ_C-P_C}{U-L}F(L)
      +\frac{P_C-LQ_C}{U-L}F(U),
\end{equation}
where $Q_C=\gamma^{\otimes2}(C)$ and
$P_C=\int_C A\,d\gamma^{\otimes2}$. The same bound holds for $G$.
For $C=(x_i,x_{i+1}]\times(x_j,x_{j+1}]$, put
\begin{align*}
 Q_i&=\Phi(x_{i+1})-\Phi(x_i),\\
 P_i&=(1-q)Q_i+q[\Phi(x_{i+1}-\mu)-\Phi(x_i-\mu)].
\end{align*}
Then $Q_C=Q_iQ_j$, $P_C=P_iP_j$, and monotonicity of $r$ gives
\[
 L=r(x_i)r(x_j),\qquad U=r(x_{i+1})r(x_{j+1}).
\]
We insert outward bounds on these quantities into~\eqref{eq:rectangle-chord}.
Outside the rectangle union, the bounds $F(A)\leq A$ and $G(A)\leq1$ control
the complement contributions.

The eleven rational edges are
\[
 -\frac{31}{10},\ \frac35,\ \frac{21}{20},\ \frac{27}{20},\
 \frac{33}{20},\ \frac{19}{10},\ \frac{43}{20},\ \frac{12}{5},\
 \frac{53}{20},\ \frac{57}{20},\ \frac{31}{10}.
\]
They form one hundred rectangles. The next subsection specifies the
interval rules; Table~\ref{tab:rational-cells} gives the resulting bounds.

\subsection{Elementary interval rules}
All interval endpoints are rational. Addition and subtraction use the usual
outward endpoint rules; multiplication takes the minimum and maximum of the
four endpoint products; reciprocation of $[l,u]$ with $l>0$ gives
$[1/u,1/l]$. The required signs are checked before reciprocation.

\paragraph{Exponential and logarithm.}
Choose integers $b\geq0$, $n\geq1$ with $|x|\leq2^b$, and put
$y=|x|/2^b$. Define
\[
 L_n(y)=\sum_{k=0}^{n-1}\frac{y^k}{k!},\qquad
 U_n(y)=L_n(y)+\frac{y^n(n+1)}{n!\,n}.
\]
For $0\leq y\leq1$,
$L_n(y)\leq e^y\leq U_n(y)$: the tail starting at order $n$ is bounded by
$y^n/n!$ times a geometric sum with ratio $1/(n+1)$. Raise the endpoints to
the power $2^b$ and, for negative $x$, reciprocate with reversed endpoints.
Monotonicity handles interval inputs. The calculation uses $(b,n)=(1,20)$
for $e^{1/4}$, $(5,20)$ for the likelihood-ratio edges, and $(6,24)$ for
Gaussian density exponentials.

A logarithm bracket $[l,u]$ for a positive interval $[a,b]$ is accepted only
if the exponential upper bound at $l$ is at most $a$ and the exponential
lower bound at $u$ is at least $b$. Monotonicity of $\log$ then proves the
enclosure. The supplement lists the individual rational brackets and Taylor
orders; floating-point logarithms do not justify any bound.

\paragraph{Square roots.}
For an integer $m>0$, let $N=10^6$ and choose the integer $j$ satisfying
$j^2\leq mN^2<(j+1)^2$. Then
$\sqrt m\in[j/N,(j+1)/N]$. Reciprocal interval arithmetic gives the required
quantities
\[
 \mu=\frac{24}{\sqrt{120}},\qquad
 \mu^{-1}=\frac5{\sqrt{120}}.
\]

\paragraph{Central Gaussian CDF values.}
Use the enclosure
\[
 n_-:=\frac{39894228040143}{10^{14}}
 <\phi(0)<
 n_+:=\frac{39894228040144}{10^{14}}.
\]
For $0\leq x\leq7/5$, put
\[
 P_0(x)=x-\frac{x^3}{6}+\frac{x^5}{40}-\frac{x^7}{336}
       +\frac{x^9}{3456}-\frac{x^{11}}{42240}+\frac{x^{13}}{599040}.
\]
Taylor's theorem for $e^{-u}$, followed by integration, gives
\[
 \left|\frac{\Phi(x)-1/2}{\phi(0)}-P_0(x)\right|
 \leq\frac{x^{15}}{2^7 7!\,15}
 \leq\frac{x^{15}}{8467200}.
\]
Here $2^7 7!\,15=9676800$; the final denominator deliberately gives a
weaker rational bound. Interval multiplication with $[n_-,n_+]$ gives the central CDF bounds.
The supplement includes the normalizer enclosure and its rational square
comparisons.

\paragraph{Gaussian tails.}
For $x>0$, define
\[
 d_{k,0}(x)=x,\qquad
 d_{k,n+1}(x)=x+\frac{k}{d_{k+1,n}(x)},\qquad
 M_n(x)=\frac1{d_{1,n}(x)}.
\]
The Mills bounds used in the calculation are
\[
 \phi(x)M_{13}(x)\leq\Phi(-x)\leq\phi(x)M_{14}(x).
\]
To verify their directions, let
\[
 I_k(x)=\int_0^\infty t^k e^{-xt-t^2/2}\,dt.
\]
Integration by parts gives $1=xI_0+I_1$ and
$kI_{k-1}=xI_k+I_{k+1}$ for $k\geq1$. Every ratio $I_k/I_{k-1}$ is positive.
Iterating these identities expresses $I_0=\Phi(-x)/\phi(x)$ as the continued
fraction above with a positive last remainder. Replacing that remainder by
zero alternates the direction at each reciprocal, giving a lower bound at
odd depth and an upper bound at even depth.

The calculation uses these tail bounds for $x>7/5$. Negative arguments are
handled by $\Phi(-x)=1-\Phi(x)$, and interval arguments by evaluating the
appropriate endpoints.

\subsection{Summing the rational bounds}
Let $x_0,\ldots,x_{10}$ be the eleven edges above in increasing order.
For each rectangle, the table gives upper-bound numerators with denominator
$10^{10}$, obtained by rounding the rational chord bound upward. The bounds
are symmetric in $i,j$, so only $i\leq j$ is shown; each off-diagonal row
is counted twice.

The outward complement bounds are
\[
 B_+=\frac{2286836197}{10^{10}},\qquad
 B_-=\frac{38961008}{10^{10}}.
\]
They round upward the rational bounds
\[
 \frac{223323847325382457311}{976562500000000000000},\qquad
 \frac{974025185261562567831}{250000000000000000000000},
\]
respectively. Adding the table entries, with off-diagonal multiplicity two,
and the complement contributions yields
\begin{align}
 H_{e^{1/4}}(P^{\otimes3}, Q^{\otimes3})
 &\leq\frac{7420364678}{10^{10}}<\frac34,
 \label{eq:rational-forward-total}\\
 H_{e^{1/4}}(Q^{\otimes3}, P^{\otimes3})
 &\leq\frac{7237647581}{10^{10}}<\frac34.
 \label{eq:rational-reverse-total}
\end{align}
Thus the sampled candidate is feasible, and $v_3\leq5/32$.

For the full-batch lower bracket, evaluate the profile at $v=4/25$. Then
$a=5/2$, and the Gaussian formula gives
\[
 G_{1/4}(5/2)=\Phi(23/20)-e^{1/4}\Phi(-27/20).
\]
The central CDF bounds and the exponential interval with $(b,n)=(1,20)$ give
\[
 \Phi(23/20)-e^{1/4}\Phi(-27/20)>\frac{19}{25}>\frac34.
\]
Hence $4/25$ is infeasible at full batch, and
\eqref{eq:calibration-transfer} gives $4/25<v_2$. Together with the sampled
bounds, this proves the privacy-variance inequalities in
Theorem~\ref{thm:strict-interior-counterexample}.

\clearpage
\begingroup\small
\begin{longtable}{@{}rrrr@{}}
\caption{Directed rectangle bounds, with denominator $10^{10}$.
Off-diagonal rows are counted twice.}
\label{tab:rational-cells}\\
\toprule $i$ & $j$ & Forward numerator & Reverse numerator \\\midrule
\endfirsthead
\toprule $i$ & $j$ & Forward numerator & Reverse numerator \\\midrule
\endhead
\midrule
\multicolumn{4}{r}{\small Continued on the next page}\\
\endfoot
\bottomrule
\endlastfoot
0 & 0 & 179309587 & 4541683687 \\
0 & 1 & 75622963 & 718523085 \\
0 & 2 & 69056184 & 283059965 \\
0 & 3 & 91045268 & 147341455 \\
0 & 4 & 94610823 & 49317043 \\
0 & 5 & 113839529 & 15318582 \\
0 & 6 & 131941862 & 0 \\
0 & 7 & 139954908 & 0 \\
0 & 8 & 111297082 & 0 \\
0 & 9 & 125975330 & 0 \\
1 & 1 & 30525041 & 104528577 \\
1 & 2 & 27152594 & 35842913 \\
1 & 3 & 35097753 & 15069211 \\
1 & 4 & 36770950 & 3637515 \\
1 & 5 & 44506072 & 0 \\
1 & 6 & 49279269 & 0 \\
1 & 7 & 50820946 & 0 \\
1 & 8 & 38394831 & 0 \\
1 & 9 & 42280938 & 0 \\
2 & 2 & 23806351 & 9614344 \\
2 & 3 & 31009020 & 3290409 \\
2 & 4 & 32740303 & 0 \\
2 & 5 & 37920442 & 0 \\
2 & 6 & 41171906 & 0 \\
2 & 7 & 40662238 & 0 \\
2 & 8 & 30071802 & 0 \\
2 & 9 & 32717715 & 0 \\
3 & 3 & 40879835 & 59609 \\
3 & 4 & 40928790 & 0 \\
3 & 5 & 46327961 & 0 \\
3 & 6 & 48166599 & 0 \\
3 & 7 & 46477322 & 0 \\
3 & 8 & 33960816 & 0 \\
3 & 9 & 36690918 & 0 \\
4 & 4 & 40398228 & 0 \\
4 & 5 & 43807871 & 0 \\
4 & 6 & 44478157 & 0 \\
4 & 7 & 42346729 & 0 \\
4 & 8 & 30722365 & 0 \\
4 & 9 & 33052386 & 0 \\
5 & 5 & 46472095 & 0 \\
5 & 6 & 46582913 & 0 \\
5 & 7 & 44021815 & 0 \\
5 & 8 & 31809211 & 0 \\
5 & 9 & 34139592 & 0 \\
6 & 6 & 46337541 & 0 \\
6 & 7 & 43592151 & 0 \\
6 & 8 & 31420920 & 0 \\
6 & 9 & 33672920 & 0 \\
7 & 7 & 40898832 & 0 \\
7 & 8 & 29435883 & 0 \\
7 & 9 & 31517496 & 0 \\
8 & 8 & 21168396 & 0 \\
8 & 9 & 22654198 & 0 \\
9 & 9 & 24237093 & 0 \\
\end{longtable}
\endgroup

\begin{thebibliography}{9}
\bibitem{raisa2024}
O.~R\"ais\"a, J.~J\"alk\"o, and A.~Honkela.
\newblock Subsampling is not Magic: Why Large Batch Sizes Work for
Differentially Private Stochastic Optimisation.
\newblock In \emph{Proceedings of the 41st International Conference on
Machine Learning}, PMLR 235:41959--41981, 2024.
\newblock \url{https://proceedings.mlr.press/v235/raisa24a.html}.


\bibitem{ballewang2018}
B.~Balle and Y.-X.~Wang.
\newblock Improving the Gaussian Mechanism for Differential Privacy:
Analytical Calibration and Optimal Denoising.
\newblock In \emph{Proceedings of the 35th International Conference on
Machine Learning}, PMLR 80:394--403, 2018.
\newblock \url{https://proceedings.mlr.press/v80/balle18a.html}.


\bibitem{rational-supplement}
Rational bounds for the fixed-participation counterexample, source supplement r05.
\newblock Accompanying file:
\href{poisson-fixed-epoch-large-batch-optimality-review-data-r05.zip}{%
\nolinkurl{poisson-fixed-epoch-large-batch-optimality-review-data-r05.zip}}.

\end{thebibliography}
\end{document}
